% !TeX program = lualatex
\documentclass[aspectratio=169,10pt,spanish]{beamer}

\usepackage[spanish,es-nodecimaldot]{babel}
\usepackage{graphicx}
\usepackage{xcolor}
\usepackage{booktabs}
\usepackage{tikz}
\usepackage{hyperref}
\usepackage{array}
\usepackage{colortbl}
\usepackage{listings}
\usepackage{calc}
\usepackage{amsmath}
\usepackage{pifont}
\usepackage{fontspec}

\usetheme{Madrid}
\usecolortheme{seahorse}

% ============== CONFIGURACIÓN ==============
\definecolor{verdeobs}{HTML}{14B8A6}
\definecolor{violetaobs}{HTML}{7C3AED}
\definecolor{grisoscuro}{HTML}{1E293B}
\definecolor{grismedio}{HTML}{64748B}

\newcommand{\ok}{\textcolor{verdeobs}{\textbf{\ding{51}}}}
\newcommand{\no}{\textcolor{red}{\textbf{\ding{55}}}}
\newcommand{\flecha}{\textcolor{violetaobs}{\textbf{$\blacktriangleright$}}}

\setbeamercolor{alerted text}{fg=violetaobs}
\setbeamercolor{frametitle}{bg=grisoscuro,fg=white}
\setbeamercolor{progress bar}{fg=verdeobs}

\title{\textbf{Cross-Bright Consolidación PTE}}
\subtitle{Normalización Presupuesto Sector Educación 2015-2024}
\author{Observatorio de Datos de Educación en Colombia}
\date{\today}
\institute{}

% ============== INICIO DOCUMENTO ==============
\begin{document}

\maketitle

% ============== DIAPOSITIVA 1: PROBLEMA ==============
\begin{frame}{Problema Original}
\begin{columns}[T]
\column{0.5\textwidth}
\textbf{Antes:}
\begin{itemize}
  \item 2015-2018: Archivos Excel manuales del MEN
  \item Encoding latin1, nombres de columnas desordenados
  \item 2019-2024: API oficial SIIF, esquema completamente distinto
  \item Cada año tenia sus propios campos, sin estandarización
  \item No habia forma de unir ni comparar historico
\end{itemize}
\column{0.5\textwidth}
\textbf{Ahora:}
\begin{itemize}
  \item \ok Esquema unificado único para los 10 años
  \item \ok 43.074 registros oficiales reales
  \item \ok 0 años nulos
  \item \ok Total pagos: \textbf{2.815,1 Billones COP}
\end{itemize}
\end{columns}
\end{frame}

% ============== DIAPOSITIVA 2: DRIFT ESQUEMA ==============
\begin{frame}{Drift de Esquema Histórico}
Esta es la razón por la que existía el problema:
\begin{center}
\includegraphics[width=0.95\textwidth]{../Imagenes/Esquema_Drift.png}
\end{center}
\vspace{0.2cm}
\small
✅ Se ve claramente el corte exacto en 2019 donde el MEN cambió de sistema, desaparecieron 5 campos antiguos y aparecieron 8 campos nuevos. Cross-Bright soluciona exactamente esto.
\end{frame}

% ============== DIAPOSITIVA 3: CROSS BRIGHT ==============
\begin{frame}{Metodología Cross-Bright}
Es el algoritmo de normalización implementado en \texttt{header\_utils.py}:

\vspace{0.5cm}
\begin{tabular}{@{}lp{10cm}@{}}
\textbf{1. Slugify} & Minúsculas, sin espacios, sin tildes, ASCII limpio \\
\textbf{2. Mapeo} & Diccionario de +200 variantes históricas → nombre canónico \\
\textbf{3. Preservación} & Columnas no mapeadas se mantienen con nombre original \\
\textbf{4. Cero Magic} & Todo explícito, sin inferencias automaticas \\
\textbf{5. Idempotente} & Aplicarlo múltiples veces no cambia el resultado \\
\end{tabular}

\vspace{0.5cm}
\begin{block}{}
\centering \Large ✅ \textbf{NO USAMOS JOIN. USAMOS ALINEACIÓN DE ESQUEMAS.}
\end{block}
\end{frame}

% ============== DIAPOSITIVA 4: VARIABLES ESTANDARIZADAS ==============
\begin{frame}{Variables Estandarizadas al Formato Moderno SIIF}
\resizebox{\textwidth}{!}{
\begin{tabular}{lcc}
\toprule
\textbf{Nombre Canónico Final} & \textbf{Variantes históricas normalizadas} & \textbf{Cobertura} \\
\midrule
\texttt{anio\_proceso} & \texttt{anio, ano, year, anio\_reporte, vigencia} & 100\% \\
\texttt{codigo\_uej} & \texttt{uej, uej\_codigo, codigoentidad, codigo\_entidad} & 100\% \\
\texttt{nombre\_uej} & \texttt{nombreentidad, nombre\_entidad} & 100\% \\
\texttt{apropiacioninicial} & \texttt{apr\_inicial, apropiacion\_inicial} & 100\% \\
\texttt{apropiacionvigente} & \texttt{apr\_vigente, apropiacion\_vigente} & 100\% \\
\texttt{compromisos} & \texttt{compromiso} & 100\% \\
\texttt{obligaciones} & \texttt{obligacion} & 100\% \\
\texttt{pagos} & \texttt{pago} & 100\% \\
\texttt{fuente} & \texttt{fuentefinanciacion, fuente\_financiacion} & 100\% \\
\midrule
\texttt{sector} & (solo disponible en API 2019+) & 64.1\% \\
\texttt{tipo\_gasto} & (solo disponible en API 2019+) & 64.1\% \\
\texttt{detalle\_gasto} & (solo disponible en API 2019+) & 64.1\% \\
\midrule
\texttt{orden\_pago} & (solo disponible en Excel 2015-2018) & 7.7\% \\
\texttt{cdp} & (solo disponible en Excel 2015-2018) & 7.7\% \\
\texttt{apropiacionbloqueada} & (solo disponible en Excel 2015-2018) & 7.7\% \\
\bottomrule
\end{tabular}
}
\end{frame}

% ============== DIAPOSITIVA 5: FLUJO ==============
\begin{frame}{Proceso de Consolidación Paso a Paso}
\begin{center}
\begin{tikzpicture}[node distance=1.2cm, auto, >=stealth,
  block/.style={rectangle, draw, thick, fill=grisoscuro!10, rounded corners, minimum height=1cm, align=center}]

\node[block] (excel) {Raw Excel \\ 2015-2018};
\node[block] (api) [right=of excel] {API JSON \\ 2019-2024};
\node[block] (mapper) [below=of $(excel)!0.5!(api)$] {Cross-Bright Mapper};
\node[block] (encoding) [below=of mapper] {Normalizador Encoding};
\node[block] (validador) [below=of encoding] {Validador Tipos Numéricos};
\node[block] (dedupe) [below=of validador] {Deduplicador Natural};
\node[block] (final) [below=of dedupe] {\textbf{pte\_consolidado\_full.parquet}};

\draw[->, thick] (excel) -- (mapper);
\draw[->, thick] (api) -- (mapper);
\draw[->, thick] (mapper) -- (encoding);
\draw[->, thick] (encoding) -- (validador);
\draw[->, thick] (validador) -- (dedupe);
\draw[->, thick] (dedupe) -- (final);
\end{tikzpicture}
\end{center}
\end{frame}

% ============== DIAPOSITIVA 6: CURVA LORENZ ==============
\begin{frame}{Concentración del Gasto por Entidad}
\begin{columns}[T]
\column{0.5\textwidth}
\includegraphics[width=\textwidth]{../Imagenes/Curva_Lorenz.png}
\column{0.5\textwidth}
\vspace{1cm}
\textbf{Coeficiente Gini = 0.950}

\vspace{0.5cm}
\begin{itemize}
  \item El 5\% de las entidades concentran el 95\% del presupuesto
  \item El MEN absorbe el 62\% del total
  \item Las 4 universidades mas grandes se reparten otro 23\%
  \item Las restantes 1.240 entidades se reparten el 15\% restante
\end{itemize}
\end{columns}
\end{frame}

% ============== DIAPOSITIVA 7: TOP ENTIDADES SIN MEN ==============
\begin{frame}{Top Entidades Excluyendo MEN}
\begin{center}
\includegraphics[width=0.9\textwidth]{../Imagenes/Top_Entidades_Sin_MEN.png}
\end{center}
\end{frame}

% ============== DIAPOSITIVA 8: EVOLUCION TEMPORAL ==============
\begin{frame}{Evolución Presupuestal Histórica}
\begin{center}
\includegraphics[width=0.9\textwidth]{../Imagenes/PTE_Evolucion_Presupuestal.png}
\end{center}
\end{frame}

% ============== DIAPOSITIVA 9: EFICIENCIA ==============
\begin{frame}{Eficiencia de Ejecución}
\begin{center}
\includegraphics[width=0.9\textwidth]{../Imagenes/PTE_Eficiencia_Ejecucion.png}
\end{center}
\end{frame}

% ============== DIAPOSITIVA 10: CALIDAD ==============
\begin{frame}{Calidad de Datos}
\vspace{0.3cm}
\begin{tabular}{lcc}
\toprule
\textbf{Categoria} & \textbf{Estado} & \textbf{Observaciones} \\
\midrule
\ok Columnas 100\% completas & 11 campos & Todos los numéricos y llaves \\
\flecha Columnas solo API 2019+ & 8 campos & 64.1\% nulos \\
\flecha Columnas solo Excel 2015-2018 & 5 campos & 7.7\% nulos \\
\no Columnas 100\% nulas & 0 & --- \\
\ok Años nulos & 0 & --- \\
\ok Montos nulos & 0 & --- \\
\bottomrule
\end{tabular}

\vspace{0.5cm}
\begin{block}{Nota Importante}
Los nulos NO son fallos. Son la consecuencia real de que el MEN cambió de sistema en 2019 y eliminó y agregó campos. Esto es normal en datos gubernamentales multianuales.
\end{block}
\end{frame}

% ============== DIAPOSITIVA 11: PRÓXIMOS PASOS ==============
\begin{frame}{Próximos Pasos Recomendados}
\begin{enumerate}
  \item \flecha \textbf{Cruzar con SNIES}: Unir PTE con matriculas y graduados para medir eficiencia
  \item \flecha \textbf{Modelo Estrella}: Ya tenemos \texttt{fact\_presupuesto + dimensiones} listo
  \item \flecha Dashboard Streamlit: Implementar filtros interactivos por entidad, año, rubro
  \item \flecha Análisis causal: Relación entre presupuesto y resultados académicos Saber
\end{enumerate}

\vspace{1cm}
\begin{center}
\Huge \textbf{FIN}
\end{center}
\end{frame}

\end{document}